
\section{Application to smooth group preschemes not necessarily affine}
\SGARefTarget{sga3:xii:section:7}

\paragraph{\textbf{Theorem 7.1.}}
\SGARefTarget{sga3:xii:7.1}
\emph{Let \(S\) be a prescheme and let \(G\) be an \(S\)-group
prescheme that is smooth, separated, and of finite type over \(S\).
Suppose that \(G\) admits a maximal torus locally for the faithfully
flat quasi-compact topology. Then:}

\begin{enumerate}
\item[\emph{\SGAInlineRefTarget{sga3:xii:7.1:item:a}(a)}] \emph{The map}
\[
  T\longmapsto \operatorname{Centr}_G(T)
\]
\emph{induces a bijection from the set of maximal tori of \(G\)
onto the set of Cartan subgroups of \(G\). If \(C\) corresponds to
\(T\), then \(T\) is the unique maximal torus of \(C\).}

\item[\emph{\SGAInlineRefTarget{sga3:xii:7.1:item:b}(b)}] \emph{Let \(T,T'\) be two maximal tori of \(G\),
and let \(C,C'\) be the corresponding Cartan subgroups. Then}
\[
  \operatorname{Transp}_G(T,T')
  =\operatorname{Transp}_G(C,C')
  =\operatorname{Transp}_G(T,C').
\]
\emph{The first two terms are also identical to the strict
transporters. Moreover, the functor in question is representable by
a closed subprescheme of \(G\), smooth over \(S\). The tori \(T,T'\)
and the Cartan subgroups \(C,C'\) are locally conjugate for the étale
topology.}

\item[\emph{\SGAInlineRefTarget{sga3:xii:7.1:item:c}(c)}] \emph{A maximal torus of \(G\), and a Cartan
subgroup of \(G\), exist locally for the étale topology.}

\item[\emph{\SGAInlineRefTarget{sga3:xii:theorem:7-1:item:d}\SGAInlineRefTarget{sga3:xii:7.1:item:d}(d)}] \emph{Suppose that every finite subset of a fiber
\(G_s\) is contained in an affine open of \(G\)---for example, that
\(G\) is quasi-projective over \(S\), or that \(S\) is artinian.
Then the functor}
\[
  \mathcal T:(\mathrm{Sch})^\circ/S\longrightarrow(\mathrm{Sets})
\]
\emph{defined in \SGARefLink{sga3:xii:1.10}{1.10}, the functor of maximal tori of \(G\), is
isomorphic to the functor}
\[
  \mathcal C:(\mathrm{Sch})^\circ/S\longrightarrow(\mathrm{Sets})
\]
\emph{of Cartan subgroups of \(G\), and is representable by a smooth,
separated prescheme of finite type over \(S\). This prescheme is
quasi-projective over \(S\) when \(G\) is, and is affine over \(S\)
when \(G\) is affine over \(S\), or when \(S\) is artinian.}

\item[\emph{\SGAInlineRefTarget{sga3:xii:7.1:item:e}(e)}] \emph{Let \(u:G\to G'\) be a morphism of
\(S\)-group preschemes, where \(G'\) is smooth, separated, and of
finite type over \(S\), and suppose that}
\[
  u_s(G_s)=G'_s
\]
\emph{for every \(s\in S\), that is, that \(u_s\) is faithfully
flat. For every maximal torus \(T\) of \(G\),
\(u(T)\) is a maximal torus of \(G'\). If \(G'\) has connected
fibers, then for every Cartan subgroup \(C\) of \(G\), \(u(C)\) is a
Cartan subgroup of \(G'\); and if \(C\) is the centralizer of \(T\),
then \(u(C)\) is the centralizer of \(u(T)\). In either case the
induced morphism \(T\to u(T)\), respectively \(C\to u(C)\), is
faithfully flat.}

\item[\emph{\SGAInlineRefTarget{sga3:xii:7.1:item:f}(f)}] \emph{Under the hypotheses of \textup{\SGARefLink{sga3:xii:7.1:item:e}{(e)}}, suppose
in addition that \(\ker u\) is a central subgroup of \(G\). Then}
\[
  T\longmapsto u(T)
\]
\emph{is a bijection from the set of maximal tori of \(G\) onto the
set of maximal tori of \(G'\). If \(G'\) has connected fibers, then}
\[
  C\longmapsto u(C)
\]
\emph{is likewise a bijection from the set of Cartan subgroups of
\(G\) onto the set of Cartan subgroups of \(G'\); if \(C'=u(C)\),
then \(C=u^{-1}(C')\).}
\end{enumerate}

\paragraph{\textbf{Remarks 7.2.}}
\SGARefTarget{sga3:xii:7.2}
We shall see in Exposé XV that the conclusion of
\SGARefLink{sga3:xii:7.1}{Theorem~7.1}\textup{\SGARefLink{sga3:xii:7.1:item:d}{(d)}} remains valid without the restrictive
condition on the finite subsets of the \(G_s\). We shall also prove
there the conclusions concerning Cartan subgroups alone that occur
in \textup{\SGARefLink{sga3:xii:7.1:item:b}{(b)}}, \textup{\SGARefLink{sga3:xii:7.1:item:c}{(c)}}, \textup{\SGARefLink{sga3:xii:7.1:item:d}{(d)}}, and \textup{\SGARefLink{sga3:xii:7.1:item:e}{(e)}}, under
the weaker hypothesis that \(G\) admit a Cartan subgroup locally for
the faithfully flat quasi-compact topology, but not necessarily a
maximal torus. For that purpose we shall moreover need to use \SGARefLink{sga3:xii:7.1}{7.1}
when \(S\) is artinian. The method of Exposé XV also considerably
simplifies the proof of \textup{\SGARefLink{sga3:xii:7.1:item:c}{(c)}} and \textup{\SGARefLink{sga3:xii:7.1:item:d}{(d)}} when \(S\) is
not artinian.

\begin{proof}[Proof of \SGARefLink{sga3:xii:7.1}{Theorem 7.1}]
\emph{\SGAInlineRefTarget{sga3:xii:7.2:item:a}(a)} Proceeding as in \SGARefLink{sga3:xii:3.2}{3.2}, one sees that for every maximal
torus \(T\) of \(G\),
\[
  C=\operatorname{Centr}_G(T)
\]
is indeed a Cartan subgroup of \(G\), and that \(T\) is determined
by \(C\) as the unique maximal torus of \(C\). It remains only to
show that every Cartan subgroup \(C\) of \(G\) is defined by a
maximal torus of \(G\), or equivalently that \(C\) admits a central
maximal torus.

The question is still local for the fpqc topology, so we may suppose
that \(G\) admits a maximal torus \(T\). By \SGARefLink{sga3:xi:theorem:6:2}{XI~6.2},
\[
  \operatorname{Transp}_G(T,C)
\]
is representable by a closed subprescheme \(S'\) of \(G\), smooth
over \(S\); moreover \(S'\to S\) is surjective by \SGARefLink{sga3:xii:6.6:item:c}{6.6(c)}. After the
base change \(S'\to S\), we may therefore suppose that there is a
section \(g\) of \(G\) such that
\[
  \operatorname{ad}(g)\cdot T\subset C.
\]
Then \(\operatorname{ad}(g)\cdot T\) is a maximal torus of \(C\).
It is central fiber by fiber by \SGARefLink{sga3:xii:6.6:item:c}{6.6(c)}, and hence central by
\SGARefLink{sga3:ix:theorem:5:6:item:b}{IX~5.6(b)}. This proves \textup{\SGARefLink{sga3:xii:7.2:item:a}{(a)}}.

\emph{\SGAInlineRefTarget{sga3:xii:7.2:item:b}(b)} Let \(T,T'\) be two maximal tori of \(G\), and \(C,C'\)
their corresponding Cartan subgroups. The following conditions are
equivalent:
\[
\begin{array}{ccccc}
(1)\ T\subset T',&
(2)\ T=T',&
(3)\ T\subset C',&
(4)\ C\subset C',&
(5)\ C=C'.
\end{array}
\]
This follows immediately from \textup{\SGARefLink{sga3:xii:7.2:item:a}{(a)}}. Applying the same result
after arbitrary base change gives the asserted identities among the
transporters and strict transporters.

We have already observed in \textup{\SGARefLink{sga3:xii:7.2:item:a}{(a)}} that
\(\operatorname{Transp}_G(T,C')\) is representable by a closed
subprescheme of \(G\), smooth over \(S\), whose structural morphism
is surjective. By Hensel's lemma this prescheme has a section locally
for the étale topology. Thus \(T,T'\), on the one hand, and \(C,C'\),
on the other, are locally conjugate for the étale topology. This
proves \textup{\SGARefLink{sga3:xii:7.2:item:b}{(b)}}.

\emph{\SGAInlineRefTarget{sga3:xii:7.2:item:c}(c)} Suppose first that \(S\) is artinian and local. If \(G\)
admits a maximal torus \(T\), the conjugacy theorem just proved shows
that the functor \(\mathcal T\) of maximal tori of \(G\) is
represented by the homogeneous space \(G/N\), where \(N\) is the
normalizer of \(T\) in \(G\); it is smooth by \textup{\SGARefLink{sga3:xii:7.2:item:b}{(b)}}.
Moreover, as observed in VI\(_A\)~\SGARefLink{par:sga3-VIA-3.2.1}{3.2.1}, every finite subset of
\(\mathcal T\) is contained in an affine open, because \(\mathcal T\)
is a homogeneous space under \(G\).

In general there is a finite extension \(k'/k\) of the residue field
of \(S\) such that \(G_{k'}\) has a maximal torus. The extension
\(k'/k\) is obtained by a finite flat base change \(S'\to S\), and
the maximal torus of \(G_{k'}\) lifts to a maximal torus of
\(G_{S'}\) by \SGARefLink{sga3:xi:proposition:2:1:bis}{XI~2.1 bis}. Thus \(\mathcal T_{S'}\) is represented by
a smooth, separated prescheme of finite type over \(S'\), every
finite subset of which is contained in an affine open. Its natural
descent datum is effective. Hence \(\mathcal T\) is representable,
and descent shows that it is smooth, separated, and of finite type
over \(S\). By smoothness and Hensel's lemma, \(\mathcal T\) has a
section locally for the étale topology. This proves \textup{\SGARefLink{sga3:xii:7.2:item:c}{(c)}} and
\textup{\SGARefLink{sga3:xii:7.1:item:d}{(d)}} when \(S\) is artinian. We shall prove below that in
this case \(\mathcal T\) is in fact affine over \(S\).

Now let \(S\) be arbitrary. To prove \textup{\SGARefLink{sga3:xii:7.2:item:c}{(c)}} and \textup{\SGARefLink{sga3:xii:7.1:item:d}{(d)}},
which are local assertions on \(S\) for the Zariski topology, we may
suppose that \(S\) is affine, that there is a prime \(\ell\) prime
to all residue characteristics of \(S\), and that the reductive rank
of the fibers of \(G\) is constant, say equal to \(r\). The latter
rank is visibly locally constant because of the assumed fpqc-local
existence of a maximal torus.

Choose a faithfully flat quasi-compact morphism \(S'\to S\), with
\(S'\) affine, such that \(G_{S'}\) admits a maximal torus \(T_0\),
and let \(C_0\) be its centralizer. I claim that there is a suitable
power \(n\) of \(\ell\) such that
\[
  C_0=\operatorname{Centr}_{G_{S'}}({}_nT_0).
\]
Indeed, this reduces immediately to the case where \(S'\) is
noetherian, which was treated in \SGARefLink{sga3:xi:theorem:6:2}{XI~6.2}. Fix such an \(n\), and put
\[
  M=(\mathbb Z/n\mathbb Z)^r,
  \qquad
  P=\operatorname{Hom}_{S\text{-gr}}(M_S,G).
\]
Then \(P\) is represented by a closed subprescheme of finite
presentation of \(({}_nG)^r\), where
\[
  {}_nG=\operatorname{Hom}_{S\text{-gr}}
    ((\mathbb Z/n\mathbb Z)_S,G).
\]
This is also the kernel of the \(n\)-th-power morphism on \(G\), and
is therefore represented by a closed subprescheme of finite
presentation of \(G\). Moreover, \(P\) is smooth over \(S\) by
\SGARefLink{sga3:xi:proposition:2:1}{XI~2.1}.

Let \(s\in S\). From what precedes, there is a finite separable
extension \(k''/\kappa(s)\) for which \(G_{k''}\) has a maximal
torus \(T''\). Since \({}_nT''\) is étale over \(k''\), we may,
after replacing \(k''\) by a finite separable extension, suppose
that
\[
  {}_nT''\simeq M_{k''}.
\]
Choose an étale morphism \(S''\to S\) and a point \(s''\in S''\)
above \(s\), inducing the residue-field extension
\(\kappa(s'')\simeq k''\). We obtain a section of
\(P_{S''}\otimes\kappa(s'')\) over
\(\operatorname{Spec}\kappa(s'')\). By Hensel's lemma, after
replacing the pointed prescheme \((S'',s'')\) by an étale
neighborhood if necessary, this section extends to a section of
\(P_{S''}\). Equivalently, there is a homomorphism
\[
  M_{S''}\longrightarrow G_{S''}
\]
inducing the given isomorphism
\[
  M_{\kappa(s'')}\simeq{}_nT''.
\]
By \SGARefLink{sga3:ix:theorem:6:4}{IX~6.4}---which here reduces to Nakayama's lemma---this
homomorphism is a closed immersion above an
open neighborhood of \(s''\), which we may take to be all of
\(S''\). Let \(H''\) be its image, and let \(C''\) be its
centralizer in \(G_{S''}\); by \SGARefLink{sga3:xi:theorem:6:2}{XI~6.2}, \(C''\) is a closed smooth
subprescheme of \(G_{S''}\).

\paragraph{\textbf{Lemma 7.3.}}
\SGARefTarget{sga3:xii:7.3}
\emph{Under the preceding hypotheses on \(\ell,n,r\), for every
prescheme \(S''\) over \(S\) and every maximal torus \(T''\) of
\(G_{S''}\), one has}
\[
  \operatorname{Centr}_{G_{S''}}(T'')
  =
  \operatorname{Centr}_{G_{S''}}({}_nT'').
\]

Indeed, after faithfully flat descent from \(S'\times_S S''\), we
reduce to the case in which \(G_{S''}\) admits a maximal torus
\(T''_1\) for which the displayed equality holds. By
\textup{(b)}, \(T''\) and \(T''_1\) are locally conjugate for the
étale topology, so the same equality holds for \(T''\).

Apply this result over \(\operatorname{Spec}\kappa(s'')\). It shows
that the fiber \(C''_{s''}\) is a Cartan subgroup of \(G_{s''}\).
Now, since \(G\) is smooth over \(S\), the union of the identity
components of the fibers \(G_s\) is an open subprescheme of \(G\);
this is the general EGA IV result for smooth morphisms. This open is visibly
stable under the group law and hence is a subgroup, denoted \(G^0\).
It satisfies the same preliminary hypotheses as \(G\), and maximal
tori of \(G\) correspond bijectively to maximal tori of \(G^0\).
Thus, for the proof of \textup{(c)} and \textup{(d)}, we may and do
suppose that \(G\) has connected fibers. Its Cartan subgroups then
also have connected fibers by \SGARefLink{sga3:xii:6.6:item:a}{6.6(a)}.

I claim that \(C''{}^0\) is a Cartan subgroup of \(G_{S''}\) above
an open neighborhood of \(s''\). This will complete the proof of
\textup{\SGARefLink{sga3:xii:7.2:item:c}{(c)}}. The claim follows from the next lemma.

\paragraph{\textbf{Lemma 7.4.}}
\SGARefTarget{sga3:xii:7.4}
\emph{Under the hypotheses of \SGARefLink{sga3:xii:7.1}{7.1}, suppose that \(G\) has connected
fibers. Let \(C\) be a closed group subprescheme of \(G\), smooth
over \(S\), and let \(s\in S\) be such that \(C_s\) is a Cartan
subgroup of \(G_s\). Then \(C^0\) is a Cartan subgroup of \(G\) over
an open neighborhood of \(s\).}

We reduce easily to the case where \(S\) is local and \(s\) is its
closed point, and to proving there that \(C^0\) is a Cartan subgroup
of \(G\). By flat descent we may further suppose that \(G\) admits a
maximal torus \(T\). By \SGARefLink{sga3:xi:theorem:6:2}{XI~6.2},
\[
  \operatorname{Transp}_G(T,C)
\]
is represented by a closed subprescheme of \(G\), smooth over \(S\).
The fiber of this transporter at \(s\) is nonempty by the hypothesis
on \(C_s\) and the conjugacy \SGARefLink{sga3:xii:6.6:item:a}{theorem 6.6(a)}. Faithfully flat descent
therefore reduces us to the case where the transporter has a section
over \(S\), hence to the case where \(C\) contains a maximal torus
\(T\) of \(G\). Then \(T\) is central in \(C^0\) by \SGARefLink{sga3:ix:theorem:5:6:item:a}{IX~5.6(a)}, so
\[
  C^0\subset \operatorname{Centr}_G(T).
\]
This is an inclusion of smooth group preschemes over \(S\), with
connected fibers and the same relative dimension---namely, the
dimension of their common fiber at \(s\). It is therefore an
equality, proving the lemma and completing \textup{\SGARefLink{sga3:xii:7.2:item:c}{(c)}}.

\emph{(d)} Retain the preceding notation and hypotheses on
\(\ell,n,r\), and the assumption that \(G\) has connected fibers.
Let
\[
  Q:(\mathrm{Sch})^\circ/S\longrightarrow(\mathrm{Sets})
\]
be the functor defined by
\[
  Q(S')=
  \left\{
  \begin{array}{c}
  \text{subgroups of multiplicative type in \(G_{S'}\)}\\
  \text{of type \(M=(\mathbb Z/n\mathbb Z)^r\)}
  \end{array}
  \right\}
\]
in the sense of \SGARefLink{sga3:ix:definition:1:4}{IX~1.4}. The rule
\[
  T\longmapsto{}_nT
\]
defines a morphism
\[
  \varphi:\mathcal T\longrightarrow Q,
\]
which is a monomorphism by \SGARefLink{sga3:xii:7.3}{Lemma~7.3}.

I claim that \(Q\) is representable by a separated prescheme of
finite presentation over \(S\). As noted in the proof of
\SGARefLink{sga3:xi:proposition:3:12:item:a}{XI~3.12(a)}, there is an isomorphism
\[
  Q\simeq P_0/\Gamma,
\]
where \(P_0\) is the open and closed subprescheme of
\[
  P=\operatorname{Hom}_{S\text{-gr}}(M_S,G)
\]
corresponding to monomorphisms \(M_{S'}\to G_{S'}\), as in \SGARefLink{sga3:ix:theorem:6:8}{IX~6.8},
and
\[
  \Gamma=(\operatorname{Aut}_{\mathrm{gr}}(M))_S.
\]
The hypothesis that every finite subset of a fiber \(G_s\) lies in
an affine open is stable under passage to a closed subprescheme and
under fiber products. It is therefore inherited by \(({}_nG)^r\),
then by \(P\), and then by \(P_0\). Hence \(Q\) is represented by a
prescheme of finite presentation over \(S\). The same
argument shows that \(Q\) is quasi-projective, respectively affine,
over \(S\) when \(G\) is. In every case \(Q\) is separated over
\(S\).

\paragraph{\textbf{Lemma 7.5.}}
\SGARefTarget{sga3:xii:7.5}
\emph{The preceding morphism
\(\varphi:\mathcal T\to Q\) is representable by an open immersion.}

In other words, let \(H\) be a subgroup of multiplicative type in
\(G\), of type \(M\). We must prove that there is an open subset
\(U\subset S\) with the following property: for every \(S'\to S\),
the group \(H_{S'}\) is of the form \({}_nT_{S'}\), for a suitable
maximal torus \(T_{S'}\) of \(G_{S'}\), if and only if
\(S'\to S\) factors through \(U\).

We may suppose that the nilpotent rank of the fibers of \(G\) is
constant, say equal to \(r'\); it is locally constant by
\textup{\SGARefLink{sga3:xii:7.2:item:b}{(b)}}. Put
\[
  C=\operatorname{Centr}_G(H),
\]
a closed group subprescheme smooth over \(S\) by \SGARefLink{sga3:xi:theorem:6:2}{XI~6.2}. Replacing
\(S\) by the open and closed subset where \(C\) has relative
dimension \(r'\), we may suppose that this is its relative dimension
everywhere. Then \(H={}_nT\) for a maximal torus \(T\) of \(G\) if
and only if \(C^0\) admits a central maximal torus \(T\), everywhere
of relative dimension \(r\), and \(H={}_nT\). This gives another
description of the subfunctor \(U\).

By flat descent, we may suppose that \(G\) admits a maximal torus
\(T_1\). Consider the subfunctor
\[
  R=
  \operatorname{Transp}_G
  \bigl(T_1,\operatorname{Centr}_{C^0}(C^0)\bigr)
  \subset
  \operatorname{Transp}_G(T_1,C).
\]
The transporter on the right is represented by a smooth prescheme
over \(S\) by \SGARefLink{sga3:xi:theorem:6:2}{XI~6.2}. The one on the left is represented by the
induced open subprescheme, as follows immediately from \SGARefLink{sga3:ix:theorem:5:6:item:a}{IX~5.6(a)};
in particular it too is smooth over \(S\). Its structural morphism
is therefore open. Replacing \(S\) by its image, equipped with the
induced structure, we may suppose that this morphism is surjective.
Applying \SGARefLink{sga3:xii:1.13}{1.13} to \(C^0\), we find that \(C^0\) admits a central
maximal torus \(T\), because it admits one fpqc-locally. Its relative
dimension is that of \(T_1\), namely \(r\). The remaining condition
on \(H\) is \(H={}_nT\), and by \SGARefLink{sga3:ix:corollary:2:10}{IX~2.10} this again cuts out a suitable
open and closed subset of \(S\). This proves the lemma.

\SGARefLink{sga3:xii:7.5}{Lemma~7.5} shows that \(\mathcal T\) is represented by a separated
prescheme locally of finite presentation over \(S\), and in fact of
finite presentation. For the latter assertion, one may either
inspect the proof of \SGARefLink{sga3:xii:7.5}{Lemma~7.5} to see that \(\varphi\) is a
quasi-compact open immersion, or descend fpqc-locally to the case
where \(G\) has a maximal torus \(T\), when
\[
  \mathcal T\simeq G/\operatorname{Norm}_G(T).
\]
This expression, or \SGARefLink{sga3:xi:proposition:2:1}{XI~2.1}, also shows that \(\mathcal T\) is smooth
over \(S\). If \(G\) is quasi-projective over \(S\), then so are
\(Q\) and \(\mathcal T\). If \(G\) is affine over \(S\), the
affineness of \(\mathcal T\) over \(S\) was proved in 5.4. It is
unknown here whether, without the hypothesis that \(G/S\) is
affine, \(n\) can be chosen so that the open immersion of Lemma~7.5
is also closed.

It remains to show that \(\mathcal T\) is affine over \(S\) when
\(S\) is artinian. By EGA I~6.1.7, we reduce to the case where
\(S\) is the spectrum of a field, which we may suppose algebraically
closed. In view of \textup{(f)}, proved below, it suffices to prove
the assertion for
\[
  G/\operatorname{Centr}(G).
\]
This quotient is affine by 6.1, so the preceding affine case
applies. This proves \textup{(d)}.

\emph{\SGAInlineRefTarget{sga3:xii:7.5:item:e}(e)} By IX~6.8, there is a subtorus \(T'\) of \(G'\) such
that \(u\) induces a faithfully flat morphism
\[
  T\longrightarrow T'.
\]
This characterizes \(T'\) as the subsheaf \(u(T)\) of \(G'\). Let
\(C\) and \(C'\) be the centralizers of \(T\) and \(T'\),
respectively. We prove that \(C\to C'\) is flat, and faithfully flat
when \(G'\) has connected fibers. Since \(C\) and \(C'\) are flat
and of finite presentation over \(S\), SGA~1, I~5.9 reduces us to a
base field, which we may suppose algebraically closed. Replacing
\(G,G'\) by their identity components does not change \(C^0,C'{}^0\),
so we may suppose \(G,G'\) connected. The assertion then follows
from 6.6(d), together with \textup{(a)}.

\emph{(f)} By \textup{(a)} and \textup{(e)}, it is enough to prove
the assertion about Cartan subgroups. A Cartan subgroup \(C\) of
\(G\), being the centralizer of a maximal torus, contains the center
of \(G\), and therefore contains \(\ker u\). Thus
\[
  C=u^{-1}(C'),
  \qquad C'=u(C),
\]
where \(C'\) is the Cartan subgroup of \(G'\) supplied by
\textup{\SGARefLink{sga3:xii:7.5:item:e}{(e)}}. Hence \(C\mapsto u(C)\) is injective. To prove
surjectivity, it suffices to show that \(u^{-1}(C')\) is a Cartan
subgroup of \(G\) for every Cartan subgroup \(C'\) of \(G'\).

The question is fpqc-local. We may suppose that \(G\) admits a
Cartan subgroup \(C_1\). Then \(C'_1=u(C_1)\) is a Cartan subgroup
of \(G'\), and is fpqc-locally conjugate to \(C'\) by \textup{(b)}.
Since
\[
  u^{-1}(C'_1)=C_1
\]
is a Cartan subgroup of \(G\), the same is true of \(u^{-1}(C')\).
This proves \textup{(f)} and completes the proof of the theorem.
\end{proof}

Alternatively, to prove that \(u^{-1}(C')\) is a Cartan subgroup,
observe that it is flat over \(S\), since \(u\) is; SGA~1, I~5.9
then reduces the assertion, by definition, to the case of a base
field, where 6.6(e) applies.

\paragraph{\textbf{Corollary 7.6.}}
\SGARefTarget{sga3:xii:7.6}
\emph{Under the hypotheses of 7.1\textup{(e)}, \(G'\) admits a
maximal torus locally for the étale topology, and therefore satisfies
the preliminary hypotheses imposed on \(G\). If moreover \(\ker u\)
is central, then the functors \(\mathcal T_G,\mathcal C_G\) of
maximal tori and Cartan subgroups of \(G\) are isomorphic to the
analogous functors \(\mathcal T_{G'},\mathcal C_{G'}\) for \(G'\).
Consequently, whenever these functors are representable, they are
represented by isomorphic \(S\)-preschemes.}

\paragraph{\textbf{Remark 7.7.}}
\SGARefTarget{sga3:xii:7.7}
\emph{(a)} In contrast with the case where \(G\) is affine over
\(S\)---in fact it is enough that \(G\) have affine fibers, as will
be seen in Exposé XVI---local constancy of the reductive rank of
\(G\) does not imply that \(G\) admits a maximal torus fpqc-locally.
For example, let \(S\) be the spectrum of a discrete valuation ring,
and let \(G_1,G_2\) be smooth separated group schemes of finite type
over \(S\), such that the generic fiber of \(G_1\) is an elliptic
curve and its special fiber is \(\mathbb G_m\), while the generic
fiber of \(G_2\) is \(\mathbb G_m\) and its special fiber is
\(\mathbb G_a\). Put
\[
  G=G_1\times_S G_2.
\]
Both fibers of \(G\) have reductive rank \(1\), but \(G\) plainly
does not admit a maximal torus fpqc-locally.

It is, on the other hand, very plausible that the following
condition on a smooth separated group of finite type over a
prescheme \(S\) is sufficient for the étale-local existence of a
maximal torus: the reductive rank and the abelian rank of the fibers
of \(G\) are locally constant functions.

\emph{(b)} In the proof of \SGARefLink{sga3:xii:7.1}{7.1}, particularly in \textup{\SGARefLink{sga3:xii:7.1:item:a}{(a)}}, we
used \SGARefLink{sga3:xi:theorem:6:2}{XI~6.2} in cases where \(S\) was not assumed locally noetherian.
In the instances at hand, however, the reduction to affine
noetherian \(S\) is immediate.

Here is a variant of \SGARefLink{sga3:xii:7.1:item:b}{7.1(b)}.

\paragraph{\textbf{Proposition 7.8.}}
\SGARefTarget{sga3:xii:7.8}
\emph{Let \(G\) be a smooth \(S\)-group prescheme of finite
presentation with connected fibers, let \(H\) be a smooth
\(S\)-group prescheme, and let}
\[
  i:H\longrightarrow G
\]
\emph{be a monomorphism of \(S\)-groups, making \(H\) an
\(S\)-subgroup of \(G\). For every Cartan subgroup
\(C\) of \(G\),\footnote{The proof only requires \(C\) to be a
smooth subgroup of \(G\) whose geometric fibers are connected
centralizers of subgroups of maximal tori.} the transporter
\(\operatorname{Transp}_G(C,H)\) is represented by a closed
subprescheme of \(G\), smooth over \(S\).}

The representability of the transporter by a closed subprescheme of
\(G\), of finite presentation, follows from \SGARefLink{sga3:xi:corollary:6:11}{XI~6.11}, because \(C\)
has connected fibers when \(G\) does, by \SGARefLink{sga3:xii:6.6:item:b}{6.6(b)}. To prove smoothness,
the standard procedure reduces first to affine noetherian \(S\),
then to artinian local \(S\), and finally, by descent, to an
algebraically closed residue field. In that case \(C\) has a maximal
torus \(T\), which is also a maximal torus of \(G\). Unless the
reductive and nilpotent ranks of the fiber \(H_0\) equal those of
\(G_0\), the transporter is empty. Under this equality one sees,
using the connectedness of \(C\) and the smoothness of the
centralizer in \(H\) of a maximal torus of \(H\), that
\[
  \operatorname{Transp}_G(C,H)
  =
  \operatorname{Transp}_G(T,H).
\]
The right-hand side is smooth by \SGARefLink{sga3:xi:proposition:2:5}{XI~2.5}, and therefore so is the
left-hand side.

The same argument proves the more general assertion
\SGARefLink{sga3:xii:7.9:item:b:02}{7.9(b)} below.

\paragraph{\textbf{Proposition 7.9.}}
\SGARefTarget{sga3:xii:7.9}
\emph{Let \(G\) and \(i:H\to G\) be as in \SGARefLink{sga3:xii:7.8}{7.8}. Suppose moreover
that, for every \(s\in S\), \(H_s\) is connected and has the same
reductive rank and the same nilpotent rank as \(G_s\); equivalently,
\(H_s\) contains a Cartan subgroup of \(G_s\). Then:}

\begin{enumerate}
\item[\emph{(a)}] \emph{\(\operatorname{Norm}_G(H)\) is represented
by a closed subprescheme of \(G\), of finite presentation over
\(S\), and the canonical monomorphism}
\[
  H\longrightarrow\operatorname{Norm}_G(H)
\]
\emph{is an open immersion. Consequently \(i\) is an immersion and}
\[
  H=\operatorname{Norm}_G(H)^0.
\]

\item[\emph{(b)}] \emph{For every Cartan subgroup \(C\) of \(G\),
\(\operatorname{Transp}_G(C,H)\) is a closed subprescheme of \(G\),
smooth over \(S\), with surjective structural morphism. If \(C\) is
the centralizer of a maximal torus \(T\) of \(G\), then moreover}
\[
  \operatorname{Transp}_G(T,H)
  =
  \operatorname{Transp}_G(C,H).
\]

\item[\emph{(c)}] \emph{Let \(C\) be a group subprescheme of \(H\).
Then \(C\) is a Cartan subgroup of \(H\) if and only if it is a
Cartan subgroup of \(G\).}

\item[\emph{(d)}] \emph{If \(G\) admits a Cartan subgroup,
respectively a maximal torus, locally for the étale topology or for
the fpqc topology, then so does \(H\).}
\end{enumerate}

\begin{proof}
\emph{\SGAInlineRefTarget{sga3:xii:7.9:item:a:02}(a)} The representability of \(\operatorname{Norm}_G(H)\) by
a closed subprescheme of \(G\), of finite presentation over \(S\),
is contained in \SGARefLink{sga3:xi:corollary:6:11}{XI~6.11}. Since \(H\) is smooth, hence flat and
locally of finite presentation over \(S\), VI\(_B\)~\SGARefLink{cor:sga3-VIB-2.6}{2.6} reduces the
assertion that
\[
  H\longrightarrow\operatorname{Norm}_G(H)
\]
is an open immersion to the fibers. We therefore reduce to the case
where \(S=\operatorname{Spec}k\), with \(k\) algebraically closed.

The assertion then reduces to showing that the corresponding map on
Lie algebras is an isomorphism, or equivalently that
\[
  (\mathfrak g/\mathfrak h)^H=0,
\]
where \(\mathfrak g,\mathfrak h\) are the Lie algebras of \(G,H\)
and the superscript denotes invariants; compare II~5.2.3(i). By
hypothesis \(H\) contains a Cartan subgroup \(C\) of \(G\), the
centralizer of a maximal torus \(T\) of \(G\). It therefore suffices
to prove
\[
  (\mathfrak g/\mathfrak h)^T=0.
\]
By complete reducibility of representations of \(T\)
(\SGARefLink{sga3:i:proposition:4-7-3}{I~4.7.3}), this follows from
\[
  (\mathfrak g/\mathfrak c)^T=0,
  \qquad \mathfrak c=\operatorname{Lie}(C).
\]
The latter equality is equivalent to
\[
  \mathfrak g^T=\mathfrak c.
\]
It says that the centralizer \(C\) and the normalizer of \(T\) have
the same Lie algebra, and follows from the fact that \(C\) is an
open subgroup of that normalizer by \SGARefLink{sga3:xi:proposition:5:9}{XI~5.9}. This proves
\textup{\SGARefLink{sga3:xii:7.9:item:a:02}{(a)}}.

\emph{\SGAInlineRefTarget{sga3:xii:7.9:item:b:02}(b)} was proved in \SGARefLink{sga3:xii:7.8}{7.8}.

\emph{\SGAInlineRefTarget{sga3:xii:7.9:item:c:02}(c)} Since \(H\) is a subprescheme of \(G\), the assertion
reduces immediately to an algebraically closed base field, where it
follows directly from the hypothesis on \(H\).

\emph{\SGAInlineRefTarget{sga3:xii:7.9:item:d:02}(d)} follows from \textup{\SGARefLink{sga3:xii:7.9:item:b:02}{(b)}}, \textup{\SGARefLink{sga3:xii:7.9:item:c:02}{(c)}}, and the Hensel
lemma \SGARefLink{sga3:xi:proposition:1:10}{XI~1.10}.
\end{proof}

\paragraph{\textbf{Corollary 7.10.}}
\SGARefTarget{sga3:xii:7.10}
\emph{Let \(G,H\) be as in \SGARefLink{sga3:xii:7.9}{7.9}, and suppose that for every
algebraically closed field \(k\) over \(S\), every element of
\(G_k(k)\) normalizing \(H_k\) belongs to \(H_k(k)\). Then \(H\) is
a closed subprescheme of \(G\) and is its own normalizer.}

This follows immediately from \SGARefLink{sga3:xii:7.9:item:a:02}{7.9(a)}. We shall apply it in
particular to Borel subgroups, and more generally to parabolic
subgroups, of \(G\).

\paragraph{\textbf{Corollary 7.11.}}
\SGARefTarget{sga3:xii:7.11}
\emph{Let \(G,H\) be as in \SGARefLink{sga3:xii:7.9}{7.9}. Then \(H\) contains every group
subprescheme \(Z\) of \(G\) which is central in \(G\), flat, and of
finite presentation over \(S\).}

As usual, we reduce first to affine noetherian \(S\), then to
artinian \(S\); the hypotheses of \SGARefLink{sga3:xii:7.1}{7.1} then hold. By \SGARefLink{sga3:xii:7.9:item:c:02}{7.9(c)}, we may
suppose that \(H\) contains a Cartan subgroup \(C\) of \(G\), and it
remains to show \(C\supset Z\). Since \(S\) is artinian, the quotient
\[
  G'=G/Z
\]
is represented by a group prescheme of finite type over \(S\); the
canonical morphism \(u:G\to G'\) is faithfully flat and has kernel
\(Z\) by VI\(_A\)~\SGARefLink{sga3:xii:3.2}{3.2}. The group \(G'\) is smooth over \(S\).
\SGARefLink{sga3:xii:7.1:item:f}{Theorem~7.1(f)} therefore shows that \(C=u^{-1}(C')\) for a Cartan
subgroup \(C'\) of \(G'\), and hence that \(C\) contains \(Z\).

\paragraph{\textbf{Corollary 7.12.}}
\SGARefTarget{sga3:xii:7.12}
\emph{Let \(G,G'\) be smooth \(S\)-group preschemes of finite
presentation, and let \(u:G\to G'\) be a faithfully flat group
homomorphism, meaning that \(u_s\) is faithfully flat for every
\(s\in S\). Suppose that \(\ker u\) is central and that \(G\) has
connected fibers. Then}
\[
  H'\longmapsto H=u^{-1}(H')
\]
\emph{defines a bijection between:}

\begin{itemize}
\item \emph{the group subpreschemes \(H'\) of \(G'\), smooth and of
finite presentation over \(S\), having at every \(s\in S\) the same
reductive rank and the same nilpotent rank as \(G'\); and}
\item \emph{the group subpreschemes \(H\) of \(G\), smooth and of
finite presentation over \(S\), having at every \(s\in S\) the same
reductive rank and the same nilpotent rank as \(G\).}
\end{itemize}

\emph{Moreover, \(H\) has connected fibers if and only if \(H'\)
does.}

Let \(H\) be a group subprescheme of \(G\) with the stated
properties. By \SGARefLink{sga3:xii:7.11}{7.11} it contains \(Z=\ker u\), so faithfully flat
descent gives
\[
  H=u^{-1}(H')
\]
for a uniquely determined group subprescheme \(H'\) of \(G'\).
Using \SGARefLink{sga3:xii:6.6:item:e}{6.6(e)}, one checks immediately that \(H'\) has the asserted
properties. If \(H\) has connected fibers, so does \(H'=H/Z\).

Conversely, start with \(H'\subset G'\) having the stated
properties, and put \(H=u^{-1}(H')\). In view of \SGARefLink{sga3:xii:6.6:item:e}{6.6(e)}, it remains
to show that \(H\) is smooth over \(S\), and that it has connected
fibers when \(H'\) does. Since \(u\) is flat and \(H'\) is flat over
\(S\), \(H\) is already flat over \(S\). It therefore suffices to
check smoothness, respectively connectedness, on geometric fibers,
reducing to an algebraically closed base field.

Then \(H'\) contains a Cartan subgroup \(C'\) of \(G'\), and \(H\)
contains its inverse image \(C\). By \SGARefLink{sga3:xii:6.6:item:e}{6.6(e)}, \(C\) is a Cartan
subgroup of \(G\), hence is smooth and connected. The conclusion now
follows from the next lemma.

\paragraph{\textbf{Lemma 7.13.}}
\SGARefTarget{sga3:xii:7.13}
\emph{Let \(G,G'\) be group preschemes flat and of finite
presentation over \(S\), let \(u:G\to G'\) be a faithfully flat
group homomorphism, and let \(C'\) be a group subprescheme of \(G'\),
of finite presentation over \(S\), such that
\(C=u^{-1}(C')\) is smooth over \(S\), respectively has connected
fibers. For every group subprescheme \(H'\) of \(G'\), of finite
presentation over \(S\), containing \(C'\), and smooth over \(S\),
respectively with connected fibers, its inverse image
\(H=u^{-1}(H')\) is smooth over \(S\), respectively has connected
fibers.}

As noted above, this reduces at once to the case where \(S\) is the
spectrum of a field. The group \(H\) is then a principal \(C\)-bundle
over \(H/C\), by VI\(_B\)~9, while
\[
  H/C\simeq H'/C'
\]
by Exposé IV. Since \(H'\) is smooth, respectively connected, so is
\(H'/C'\), hence so is \(H/C\). The same property holds for \(C\) by
hypothesis and therefore for the bundle \(H\). This proves the
lemma and \SGARefLink{sga3:xii:7.12}{Corollary~7.12}.

\section{Semisimple elements; unions and intersections of maximal
tori in group schemes not necessarily affine}
\SGARefTarget{sga3:xii:section:8}

Throughout this section, \(G\) denotes an \(S\)-group prescheme that
is smooth and of finite presentation over \(S\), with connected
fibers.

Suppose first that \(S\) is the spectrum of an algebraically closed
field \(k\). When \(G\) is affine, the notion of a semisimple element
of \(G(k)\) was defined in \emph{BIBLE}~4, No.~4; one checks at once
that this notion is invariant under an algebraically closed extension
\(k'/k\) of the base field. Moreover, by \emph{BIBLE}~6,
Theorem~5(c), an element \(g\in G(k)\) is semisimple if and only if
it is contained in a maximal torus of \(G\).

When \(G\) is no longer assumed affine, Chevalley's theorem gives a
canonical expression of \(G\) as an extension of an abelian variety
by a smooth connected affine algebraic group:
\begin{equation}
\tag{*}\label{sga3:xii:lemma:7-13:equation:star}
  e\longrightarrow V\longrightarrow G\longrightarrow A
  \longrightarrow e.
\end{equation}
We shall say that \(g\in G(k)\) is \emph{semisimple} if it is a
semisimple element of \(V(k)\). Since the maximal tori of \(V\) are
precisely the maximal tori of \(G\), this is equivalent to saying
that \(g\) belongs to a maximal torus of \(G\). This notion is again
invariant under algebraically closed extensions of the base field.

Suppose next that \(S=\operatorname{Spec}k\) for an arbitrary field
\(k\), and let \(g\in G\). After choosing an algebraically closed
extension \(K\) of \(\kappa(g)\), the point \(g\) is the image of a
geometric point \(g'\) of \(G\) with values in an algebraically
closed extension \(K\) of \(k\). We call \(g\) semisimple when \(g'\)
is semisimple. What precedes shows that this is independent of the
choice of \(g'\). If \(k'/k\) is any extension, the set of
semisimple elements of \(G_{k'}\) is the inverse image of the set
\(G^{\mathrm{ss}}\) of semisimple elements of \(G\).

Finally, for arbitrary \(S\), a point \(g\in G\) is called
semisimple when it is semisimple in its fiber \(G_s\). Thus, for
every morphism \(S'\to S\), the subset \(G^{\mathrm{ss}}_{S'}\) is
the inverse image of \(G^{\mathrm{ss}}\). Suppose that the functor
\(\mathcal T\) of maximal tori defined in \SGARefLink{sga3:xii:1.10}{1.10} is representable by a
prescheme of finite presentation over \(S\). This is so, for
example, if \(G\) admits a maximal torus locally for the fpqc
topology, by \SGARefLink{sga3:xii:7.1:item:d}{7.1(d)}, at least when \(G\) is quasi-projective over
\(S\). Let \(T\) be the canonical maximal torus of
\(G_{\mathcal T}\), the ``universal maximal torus of \(G\)'', and
consider the morphism
\[
  u:T\longrightarrow G
\]
induced by the projection \(G_{\mathcal T}\to G\). It follows
immediately from the definition that \(G^{\mathrm{ss}}\) is the
image of this morphism. We obtain:

\paragraph{\textbf{Proposition 8.1.}}
\SGARefTarget{sga3:xii:8.1}
\emph{The set \(G^{\mathrm{ss}}\) of semisimple elements of \(G\)
is locally constructible. It is therefore constructible when \(S\),
and hence \(G\), is quasi-compact and quasi-separated.}

As usual, we reduce to affine noetherian \(S\). By the standard
noetherian criterion for constructibility (EGA \(0_{\mathrm{III}}\),
9.2.3), we may further suppose that \(S\) is integral and need only
show that there is a nonempty open \(U\subset S\) such that
\(G^{\mathrm{ss}}|_U\) is constructible. After shrinking \(U\), and
replacing it if necessary by a finite covering, we may suppose that
\(G\) is separated over \(S\) and contains a maximal torus. By
7.1(d), \(\mathcal T\) is then represented by a prescheme of finite
presentation over \(S\); the same is true of its universal torus,
whose image in \(G\) is therefore constructible.

Suppose again that \(k\) is a field, and let \(H\) be the subgroup
of \(G\) generated by the preceding morphism (VI\(_B\)~\SGARefLink{sga3:xii:1.2}{1.2}). It is
a smooth group subscheme of \(G\), connected because the universal
torus is connected, and its formation commutes with every extension
of the base field, as does the formation of \(\mathcal T\).

When \(k\) is algebraically closed, \(H\) is also the algebraic
subgroup of \(G\) generated by the maximal tori of \(G\), or,
equivalently, by all the tori of \(G\). It is likewise the smallest
algebraic subgroup of \(G\) containing all semisimple elements of
\(G(k)\). These characterizations remain valid whenever \(k\) is
infinite, because, as will be proved in Exposé XIV, the set of
\(k\)-rational points of \(\mathcal T\) is dense in \(\mathcal T\).

The subgroup \(H\) is invariant in \(G\). Indeed, it is enough to
check this after passing to an algebraically closed field; since
\(H\) and \(G\) are then smooth, it suffices to check stability
under the inner automorphisms \(\operatorname{int}(g)\),
\(g\in G(k)\), which is immediate. One can even show that \(H\) is
characteristic in \(G\), that is, stable under
\(\operatorname{Aut}_{k\text{-}\mathrm{gr}}(G)\).

It follows from \SGARefLink{sga3:xii:6.6:item:d}{6.6(d)} that \(G/H\) has reductive rank zero. More
precisely, if \(K\) is an invariant algebraic subgroup of \(G\),
then, over an algebraically closed field, the maximal tori of \(G/K\)
are the direct images of the maximal tori of \(G\). Hence \(G/K\)
has reductive rank zero if and only if \(K\) contains all maximal
tori of \(G\), or, over an arbitrary field, if and only if
\(K\supset H\). Thus \(H\) is the smallest invariant algebraic
subgroup of \(G\) for which \(G/H\) has reductive rank zero. It is
also the smallest algebraic subgroup of \(G\) having the same
reductive rank as \(G\).

Finally, \(H\) is affine. This may again be checked after passing to
an algebraic closure. In the exact sequence \((*)\), every maximal
torus of \(G\) is contained in \(V\); consequently \(H\subset V\),
and \(V\) is affine. We have proved:

\paragraph{\textbf{Proposition 8.2.}}
\SGARefTarget{sga3:xii:8.2}
\emph{Let \(G\) be a smooth connected algebraic group over a field
\(k\), and let \(\bar k\) be an algebraic closure of \(k\). There
exists an algebraic subgroup \(H\) of \(G\) such that
\(H_{\bar k}\) is the algebraic subgroup of \(G_{\bar k}\) generated
by the maximal tori, or equivalently by all the tori, of
\(G_{\bar k}\). The group \(H\) is also characterized as the
smallest algebraic subgroup of \(G\) having the same reductive rank
as \(G\), or as the smallest invariant algebraic subgroup of \(G\)
such that \(G/H\) has reductive rank zero. It is a smooth, connected,
invariant, affine subgroup of \(G\), and its formation commutes with
every extension of the base field.}

The characterization in terms of \(G/H\) gives the following
consequences.

\paragraph{\textbf{Corollary 8.3.}}
\SGARefTarget{sga3:xii:8.3}
\emph{Let \(G\) be a smooth connected algebraic group over an
algebraically closed field \(k\). Then \(G\) has reductive rank zero
---equivalently, in the notation of \SGARefLink{sga3:xii:8.2}{8.2}, \(H=(e)\)---if and only if
\(G\) is an extension of an abelian variety by a smooth connected
unipotent algebraic group, that is, a successive extension of groups
isomorphic to the additive group \(\mathbb G_a\).}

Indeed, Chevalley's exact sequence \((*)\) reduces the assertion to
the fact that a smooth connected affine group \(V\) has reductive
rank zero if and only if it is unipotent; this is contained in
\emph{BIBLE}~\SGARefLink{sga3:xii:6.4}{6.4}, Theorem~4, \SGARefLink{sga3:xii:bibliography}{Corollary~3}. Thus, in the notation of
\SGARefLink{sga3:xii:8.2}{8.2}, \(H\) is the smallest invariant algebraic subgroup of \(G\)
such that \((G/H)_{\bar k}\) is an extension of an abelian variety
by a smooth connected affine unipotent algebraic group.

\paragraph{\textbf{Corollary 8.4.}}
\SGARefTarget{sga3:xii:8.4}
\emph{In the notation of \SGARefLink{sga3:xii:8.2}{8.2}, \(G=H\), or equivalently
\(G_{\bar k}\) is generated by its maximal tori, if and only if
\(G\) is affine and every homomorphism
\(G_{\bar k}\to\mathbb G_a\) is trivial.}

\paragraph{\textbf{Remarks 8.5.}}
\SGARefTarget{sga3:xii:8.5}
\begin{enumerate}
\item[\emph{\SGAInlineRefTarget{sga3:xii:8.5:item:a}(a)}] Let \(V\) be the largest smooth connected affine
algebraic subgroup of \(G_{\bar k}\), so that \(G_{\bar k}\) is an
extension of an abelian variety by \(V\). It is well known, by
Rosenlicht,\footnote{The editors ask that references be supplied
here.} that when \(k\) is not perfect, \(V\) is not in general
defined over \(k\): there need not exist an algebraic subgroup
\(V_0\) of \(G\) such that \((V_0)_{\bar k}=V\). If, however, \(V\)
is generated by its maximal tori---equivalently, if \(V\) has no
quotient isomorphic to \(\mathbb G_{a,\bar k}\)---then such a group
does exist, namely the group \(H\) of \SGARefLink{sga3:xii:8.2}{8.2}. Thus, in this rationality
question as in many others (compare Exposé XIV, No.~6), all the
difficulties arise from unipotent groups, hence from the additive
group, whereas the presence of sufficiently many groups of
multiplicative type makes descent work.

\item[\emph{(b)}] In the notation of \SGARefLink{sga3:xii:8.2}{8.2}, let \(H'\) be the inverse
image in \(G\) of the commutator subgroup of \(G/H\). Then \(H'\) is
the smallest invariant algebraic subgroup of \(G\) such that
\((G/H')_{\bar k}\) is a commutative extension of an abelian variety
by a smooth connected affine unipotent algebraic group. It is again
a smooth connected affine subgroup of \(G\).

Assume that the characteristic \(p\) is positive, and let \(H''\)
be the inverse image in \(G\) of \(p^n(G/H')\) for \(n\) sufficiently
large. Then \(H''/H'\) is an abelian variety, and \(H''\) is the
smallest invariant algebraic subgroup of \(G\) such that \(G/H''\)
is a commutative smooth connected affine unipotent algebraic group.
It follows that \(V=H_{\bar k}\), in the notation of \textup{\SGARefLink{sga3:xii:8.5:item:a}{(a)}},
if and only if \(H''=G\), or equivalently if and only if every
homomorphism \(G_{\bar k}\to\mathbb G_a\) is trivial.
\end{enumerate}

We now generalize to smooth groups with connected fibers the notion
of reductive center developed in \SGARefLink{sga3:xii:section:4}{Section~4}, following the idea of
\SGARefLink{sga3:xii:4.10}{4.10}. Let \(\mathcal Z\) be the subfunctor of \(G\) defined by
\[
\begin{split}
  \mathcal Z(S')=\{\, &g'\in G_{S'}(S'):\text{ for every }
  S''\to S'\text{ and every maximal torus }T''\subset G_{S''},\\
  &\text{the pullback }g''\text{ of }g'\text{ is a section of }
  T''\text{ over }S''\,\}.
\end{split}
\]
Let \(\mathcal T(S')\) be the set of maximal tori of \(G_{S'}\), and
let \(\widetilde{\mathcal T}(S')\) be the set of pairs \((T',g')\)
where \(T'\) is a maximal torus of \(G_{S'}\) and \(g'\) is a section
of \(T'\) over \(S'\). Then \(\widetilde{\mathcal T}\) is a
subfunctor of
\[
  \mathcal T\times_S G=\mathcal T_G,
\]
and in this notation one may also write
\[
  \mathcal Z=
  \prod_{\mathcal T_G/G}
  \widetilde{\mathcal T}/\mathcal T_G.
\]
Using \SGARefLink{sga3:xi:theorem:6:8}{XI~6.8} and \SGARefLink{sga3:xii:7.1:item:d}{7.1(d)}, which under suitable hypotheses represents
\(\mathcal T\) by a smooth \(S\)-prescheme, one could deduce
representability of \(\mathcal Z\) under corresponding hypotheses.
We shall instead obtain it more directly below.

\paragraph{\textbf{Definition 8.6.}}
\SGARefTarget{sga3:xii:8.6}
\emph{Let \(G\) be an \(S\)-group prescheme, smooth and of finite
presentation with connected fibers. We say that \(G\) admits a
\emph{reductive center} if the preceding functor \(\mathcal Z\),
which is evidently a subgroup of \(G\), is representable by a group
of multiplicative type. In that case \(\mathcal Z\) is called the
reductive center of \(G\).}

If \(Z\) is a reductive center of \(G\), then \(Z_{S'}\) is a
reductive center of \(G_{S'}\) for every base change \(S'\to S\).
The existence of a reductive center is consequently local for the
fpqc topology. The terminology is justified because \(\mathcal Z\)
is invariant under
\[
  \operatorname{Aut}_S(G),
\]
and hence is an invariant subgroup of \(G\); by \SGARefLink{sga3:ix:corollary:5:5}{IX~5.5}, it is
central.

\SGARefLink{sga3:xii:4.5}{Lemma~4.5} is replaced here by the following statement.

\paragraph{\textbf{Lemma 8.7.}}
\SGARefTarget{sga3:xii:8.7}
\emph{Let \(u:H\to G\) be a central homomorphism, where \(H\) is of
multiplicative type and of finite type over \(S\). Suppose that for
every algebraically closed field \(k\) over \(S\), the subgroup
\(u_k(H_k)\) is contained in the largest smooth connected affine
subgroup of \(G_k\). Then \(u\) factors through every maximal torus
\(T\) of \(G\). In particular, \(u\) factors through the subfunctor
\(\mathcal Z\) defined above.}

An easy variant of \SGARefLink{sga3:ix:theorem:5:1:bis}{IX~5.1 bis}, with equality replaced by inclusion,
reduces the assertion to the case in which \(S\) is the spectrum of
an algebraically closed field: first reduce to affine noetherian
\(S\), then to artinian \(S\), and use \SGARefLink{sga3:ix:theorem:3:6}{IX~3.6}. Since \(T\) is
contained in the largest smooth connected affine subgroup \(V\) of
\(G\), we are reduced to \(G=V\), hence to the affine case proved
in \SGARefLink{sga3:xii:4.5}{4.5}.

\paragraph{\textbf{Proposition 8.8.}}
\SGARefTarget{sga3:xii:8.8}
\emph{Let \(G\) be an \(S\)-group prescheme, smooth and of finite
presentation with connected fibers.}

\begin{enumerate}
\item[\emph{(a)}] \emph{If \(S=\operatorname{Spec}k\), then \(G\)
admits a reductive center \(Z\). When \(G\) is an extension of an
abelian variety by a smooth connected affine algebraic group \(V\)
---for example, when \(k\) is algebraically closed---then \(Z\) is
also the reductive center of \(V\), and it is the largest central
subgroup of multiplicative type of \(V\).}

\item[\emph{(b)}] \emph{Let \(Z\) be a group subprescheme of
multiplicative type in \(G\). Then \(Z\) is a reductive center of
\(G\) if and only if \(Z_s\) is a reductive center of \(G_s\) for
every \(s\in S\). It is then the largest subgroup of multiplicative
type \(K\) in \(G\) such that \(K_s\) is contained in the reductive
center of \(G_s\) for every \(s\in S\). More generally, let
\(u:H\to G\) be a homomorphism with \(H\) of multiplicative type and
of finite type over \(S\), and suppose that for every algebraically
closed field \(k\) over \(S\), \(u_k\) factors through the largest
smooth connected affine subgroup of \(G_k\). Then \(u\) factors
through \(Z\), and in particular is central.}

\item[\emph{(c)}] \emph{If \(G\) admits a maximal torus locally for
the fpqc topology, then \(G\) admits a reductive center.}

\item[\emph{(d)}] \emph{If \(T\) is a maximal torus of \(G\), then}
\[
  T\cap\operatorname{Centr}(G)
  =\ker\bigl(T\longrightarrow\operatorname{GL}(\mathfrak g)\bigr),
\]
\emph{and this group is a reductive center of \(G\).}

\item[\emph{(e)}] \emph{Let \(Z\) be a reductive center of \(G\),
and suppose that \(G'=G/Z\) is representable, for example when \(S\)
is artinian. Then the unit subgroup is a reductive center of \(G'\),
and}
\[
  T'\longmapsto T=u^{-1}(T')
\]
\emph{is a bijection from the set of maximal tori of \(G'\) to the
set of maximal tori of \(G\).}
\end{enumerate}

\begin{proof}
\emph{\SGAInlineRefTarget{sga3:xii:8.8:item:a:02}(a)} Let \(S=\operatorname{Spec}k\). To prove existence, we
may pass to an algebraic closure of \(k\), and thus suppose that
\(G\) is an extension of an abelian variety by a smooth connected
affine algebraic group \(V\). For every \(S'\) over \(k\), the
maximal tori of \(G_{S'}\) are those of \(V_{S'}\): this follows
from \SGARefLink{sga3:ix:corollary:5:2}{IX~5.2} and from the fact that a homomorphism from a torus to an
abelian variety over an algebraically closed field is trivial.
Consequently the functor \(\mathcal Z\) defined using \(G\) is the
same as the one defined using \(V\), and we reduce to affine \(G\).

The functor \(\mathcal T\) is representable and is necessarily
``essentially free'' over \(k\) (\SGARefLink{sga3:viii:definition:6:1}{VIII~6.1}); hence
\(\mathcal T_G\) is essentially free over \(G\). Since the universal
torus is a closed subscheme of
\(\mathcal T_G=G_{\mathcal T}\) (for example by \SGARefLink{sga3:viii:corollary:5:7}{VIII~5.7}),
\SGARefLink{sga3:viii:theorem:6:4}{VIII~6.4} shows that
\[
  \mathcal Z=
  \prod_{\mathcal T_G/G}
  \widetilde{\mathcal T}/\mathcal T_G
\]
is represented by a closed subscheme of \(G\). It is therefore a
group subscheme. It is of multiplicative type: after passing to an
algebraic closure, \(G\) admits a maximal torus \(T\), and by
definition \(Z\subset T\); a subgroup of a group of multiplicative
type is of multiplicative type (IX~8). Thus \(Z\) is a reductive
center. \SGARefLink{sga3:xii:8.7}{Lemma~8.7} shows that it is the largest central subgroup of
multiplicative type of \(V\).

\emph{\SGAInlineRefTarget{sga3:xii:8.8:item:b:02}(b)} Only sufficiency requires proof. Suppose that \(Z\) is
of multiplicative type and that \(Z_s\) is the reductive center of
\(G_s\) for every \(s\in S\). \SGARefLink{sga3:xii:8.7}{Lemma~8.7} first shows that \(Z\) is
contained in every maximal torus of \(G\), and this remains true
after every base change. It remains to show that if \(g\) is a
section of \(G\) over \(S\) whose pullback to every \(S'\) over
\(S\) is a section of every maximal torus of \(G_{S'}\), then \(g\)
is a section of \(Z\).

As usual, since the fpqc sheaf
\[
  \mathcal Z_0=
  \prod_{\mathcal T_G/G}
  \widetilde{\mathcal T}/\mathcal T_G
\]
commutes with filtered direct limits of rings, we reduce successively
to affine, noetherian, and artinian \(S\). By \SGARefLink{sga3:xii:7.1:item:d}{7.1(d)},
\(\mathcal T\) is then represented by a smooth prescheme over \(S\).
Arguing as in \textup{\SGARefLink{sga3:xii:8.8:item:a:02}{(a)}}, \(\mathcal Z_0\) is represented by a
closed subscheme \(Z_0\) of \(G\). We already know \(Z\subset Z_0\);
by hypothesis \(Z_k=(Z_0)_k\) over the residue field \(k\), and
\(Z\) is flat over \(S\). It follows that \(Z=Z_0\), proving that
\(Z\) is the reductive center. The remaining assertions of
\textup{\SGARefLink{sga3:xii:8.8:item:b:02}{(b)}} follow from \SGARefLink{sga3:xii:8.7}{8.7}.

\emph{\SGAInlineRefTarget{sga3:xii:8.8:item:c:02}(c)} follows immediately from \textup{\SGARefLink{sga3:xii:8.8:item:d:02}{(d)}}.

\emph{\SGAInlineRefTarget{sga3:xii:8.8:item:d:02}(d)} Here \(\mathfrak g\) denotes the Lie algebra of \(G\),
and \(T\to\operatorname{GL}(\mathfrak g)\) is induced by the adjoint
representation. Trivially,
\[
  T\cap\operatorname{Centr}(G)
  \subset
  \ker\bigl(T\longrightarrow\operatorname{GL}(\mathfrak g)\bigr).
\]
The proof of the reverse inclusion is the same as in \SGARefLink{sga3:xii:4.7:item:d}{4.7(d)}. Put
\[
  Z=\ker\bigl(T\longrightarrow\operatorname{GL}(\mathfrak g)\bigr).
\]
Then \(Z\) is of multiplicative type (\SGARefLink{sga3:ix:theorem:6:8}{IX~6.8}). By \textup{\SGARefLink{sga3:xii:8.8:item:b:02}{(b)}}, we
may check that it is a reductive center over a field. Let \(Z_0\)
be the reductive center supplied by \textup{\SGARefLink{sga3:xii:8.8:item:a:02}{(a)}}. Plainly
\[
  Z_0\subset T\cap\operatorname{Centr}(G)=Z.
\]
Conversely, \(Z\) is a central subgroup of multiplicative type
contained in the smooth connected affine group \(T\), so
\textup{\SGARefLink{sga3:xii:8.8:item:a:02}{(a)}} gives \(Z\subset Z_0\). Thus \(Z=Z_0\).

\emph{\SGAInlineRefTarget{sga3:xii:8.8:item:e:02}(e)} Under the hypotheses of \SGARefLink{sga3:xii:7.1:item:f}{7.1(f)}, put \(Z=\ker u\), and
suppose that for every algebraically closed field \(k\) over \(S\),
\(Z_k\) is contained in a maximal torus of \(G_k\). One checks
readily that
\[
  T'\longmapsto u^{-1}(T')
\]
is a bijection from the maximal tori of \(G'\) to those of \(G\).
Applying this observation to the situation of \textup{\SGARefLink{sga3:xii:8.8:item:e:02}{(e)}} proves
the assertion.
\end{proof}

\paragraph{\textbf{Remarks 8.9.}}
\SGARefTarget{sga3:xii:8.9}
\begin{enumerate}
\item[\emph{(a)}] The proof of \SGARefLink{sga3:xii:8.8}{8.8} is independent of the results of
Section~4; in particular, the proof of \SGARefLink{sga3:xii:8.8:item:a:02}{8.8(a)} does not use \SGARefLink{sga3:xii:4.4}{4.4},
whose proof is somewhat laborious.

\item[\emph{(b)}] The subfunctor \(\mathcal Z\) considered in \SGARefLink{sga3:xii:8.6}{8.6}
is always central, whether or not it is representable.\ It may
therefore be tempting to call it the reductive center of \(G\) in
all cases.

\item[\emph{(c)}] One may likewise generalize \SGARefLink{sga3:xii:4.9}{4.9}. If \(G\) is a
smooth connected algebraic group over an algebraically closed
field, then \(G\) is an extension of an abelian variety by a smooth
connected unipotent algebraic group---equivalently, \(G\) has
reductive rank zero, by \SGARefLink{sga3:xii:8.3}{8.3}---if and only if the reductive center of
\(G\) is the unit group and the Lie algebra of \(G\) is nilpotent.
\end{enumerate}

\section{Complement: actions of group schemes and fixed points}
\SGARefTarget{sga3:xii:section:9}

The purpose of this section, added in January 2008, is to study fixed
points under the action of a group scheme.

\subsection{Representability of the fixed-point functor}
\SGARefTarget{sga3:xii:9.1}

Let \(S\) be a scheme, let \(G\) be an \(S\)-group scheme, and
suppose that \(G\) acts on an \(S\)-scheme \(X\). Define the
subfunctor \(X^G\) of fixed points of \(X\) under \(G\) by requiring,
for every \(S\)-scheme \(T\), that \(X^G(T)\) be the subset of
\(X(T)\) consisting of the points fixed under \(G\) (VIII~6\SGARefLink{sga3:xii:8.8:item:e:02}{(e)}).

Recall the notion of an \(S\)-pure scheme introduced in
Raynaud--Gruson, \S3.3, and in Raynaud. Let \(G\) be a flat
\(S\)-group scheme of finite presentation. Since the fibers of \(G\)
have no embedded components, purity takes the following form:
\(G\) is \emph{\(S\)-pure} if, for every strictly henselian local
\(S\)-scheme \(T\) with closed point \(t\), every generic point
\(\widetilde x\) of a fiber of \(G\times_S T\) specializes to a
point of \(G_t\); equivalently, the closure of \(\widetilde x\) in
\(G\times_S T\) meets \(G_t\).

\paragraph{\textbf{Examples.}}
\begin{enumerate}
\item[\textup{(1)}] If \(G\) is quasi-finite and separated over
\(S\), then \(G\) is \(S\)-pure if and only if it is finite over
\(S\).

\item[\textup{(2)}] If \(S=\operatorname{Spec}R\) and
\(G=\operatorname{Spec}A\), then \(G\) is \(S\)-pure if and only if
\(A\) is a projective \(R\)-module (Raynaud--Gruson, 3.3.5). In
particular, a diagonalizable \(G\) is \(S\)-pure.

\item[\textup{(3)}] Every group scheme of multiplicative type is
\(S\)-pure, since purity is local for the étale topology on \(S\).

\item[\textup{(4)}] The scheme \(G\) is \(S\)-pure if it is proper
over \(S\), if its fibers are irreducible, or if \(S\) is
semilocal artinian.

\item[\textup{(5)}] In particular, every reductive \(S\)-group
scheme is \(S\)-pure.
\end{enumerate}

\paragraph{\textbf{Proposition 9.2.}}
\SGARefTarget{sga3:xii:9.2}
\emph{Let \(G\) be an \(S\)-group scheme of finite presentation,
flat and \(S\)-pure over \(S\), acting on an \(S\)-scheme \(X\) of
finite presentation. Then the fixed-point functor \(X^G\) is
represented by a closed subscheme of \(X\), of finite presentation
over \(S\).}

Indeed, let \(a\in X(S)\), and let \(H\) be the group subscheme of
\(G\) fixing \(a\). The point \(a\) is fixed under \(G\) if and only
if \(H=G\). The coincidence functor of \(H\) and \(G\)---whose value
on an \(S\)-scheme \(T\) is a singleton when
\[
  H\times_S T\longrightarrow G\times_S T
\]
is bijective and is empty otherwise---is represented by a
closed subscheme of \(S\), defined by a finite-type sheaf of ideals
(Raynaud--Gruson, 4.1.1). Apply this result with \(S=X\) and with
\(a\) the universal point of \(X\), that is, the identity element
of
\[
  X(X)=\operatorname{Hom}_S(X,X).
\]

\setcounter{subsection}{2}
\subsection{The infinitesimal obstruction}
\SGARefTarget{sga3:xii:9.3}

Suppose that \(G\) is a flat \(S\)-group scheme of finite
presentation acting on the left on a smooth \(S\)-prescheme \(X\).
We shall study the formal smoothness of \(X^G\).

Let \(S_0\) be a closed subscheme of \(S\), defined by a
quasi-coherent sheaf of ideals \(\mathcal I\) with
\(\mathcal I^2=0\), and write
\[
  j:S_0\longrightarrow S,\qquad
  G_0=G\times_S S_0,\qquad
  X_0=X\times_S S_0.
\]
Let \(\epsilon_0\in X^G(S_0)\) be a point that lifts to \(X(S)\).
We shall study the obstruction to lifting \(\epsilon_0\) to
\(X^G(S)\).

Since \(X\) is smooth, the set of lifts of \(\epsilon_0\) to \(X(S)\)
is a trivial principal homogeneous space under the abelian group
\[
  \operatorname{Hom}_{\mathcal O_{S_0}}
  \bigl(\epsilon_0^*\Omega^1_{X_0/S_0},\mathcal I\bigr)
\]
(\SGARefLink{sga3:iii:proposition:0-2}{III~0.2}, \SGARefLink{sga3:iii:remark:0-3}{0.3}). Put
\[
  L_0=
  \operatorname{Hom}_{\mathcal O_{S_0}}
  \bigl(\epsilon_0^*\Omega^1_{X_0/S_0},\mathcal I\bigr).
\]
This is an \(\mathcal O_{S_0}\)-module. Since \(\epsilon_0\) is
fixed under \(G\), it is naturally a
\(G_0\)-\(\mathcal O_{S_0}\)-module. Let
\[
  \rho_0:G_0\longrightarrow\operatorname{Aut}_{S_0}(L_0),
  \qquad
  \rho:G\longrightarrow\operatorname{Aut}_S(j_*L_0)
\]
be the associated representations.

\paragraph{\textbf{Lemma 9.4.}}
\SGARefTarget{sga3:xii:9.4}
\emph{There is a canonically defined class}
\[
  c(\epsilon_0)\in H^1(G,j_*L_0)
  \simeq H^1(G_0,L_0)
\]
\emph{such that \(\epsilon_0\) lifts to \(X^G(S)\) if and only if
\(c(\epsilon_0)=0\).}

The displayed adjunction is the one in \SGARefLink{sga3:iii:lemma:1-1-2}{III~1.1.2}. We work on the
small fppf site of \(S\). For every \(S\)-scheme \(T\), flat and of
finite presentation, write
\[
  G_T=G\times_S T,\qquad X_T=X\times_S T.
\]
Consider the sheaf \(\mathcal A\) on this site defined by
\[
  \mathcal A(T)=
  \{\,\text{lifts of }\epsilon_0\times_S T\text{ to }X(T)\,\}.
\]
Since the formation of \(L_0\) commutes with flat base change,
\(\mathcal A(T)\) is a trivial principal homogeneous space under
\[
  H^0(T,j_*L_0).
\]
Moreover, because \(\epsilon_0\) is fixed, every \(g\in G(T)\) acts
by affine automorphisms on \(\mathcal A(T)\), compatibly with the
action of \(G\) on \(j_*L_0\):
\[
  g(a_T+v)=g(a_T)+\rho(g)(v),
  \qquad v\in H^0(T,j_*L_0).
\]
Since \(G\) is flat and of finite type, we may take \(T=G\) and
\(g\) equal to the universal point \(\operatorname{id}_G\). We
thereby obtain an action of the fppf sheaf \(G\) on \(\mathcal A\).

Let \(a\) be a lift of \(\epsilon_0\) to \(X(S)\), and let
\(g^\sharp\) denote the universal point of \(G\). Define
\[
  v^\sharp\in H^0(G,j_*L_0)
\]
by the relation
\[
  \rho(g^\sharp)(v^\sharp)=g^\sharp\cdot a_G-a_G.
\]
For every \(S\)-prescheme \(Y\), set
\[
  z(Y):G(Y)\longrightarrow H^0(Y,j_*L_0),
  \qquad g\longmapsto v^\sharp(g).
\]
This defines a \(1\)-cocycle \(z\in Z^1(G,j_*L_0)\) (Exposé I,
\S5). Its class \(c(\epsilon_0)\) does not depend on the choice of
\(a\). In particular, if \(\epsilon_0\) lifts to \(X^G(S)\), then
\(c(\epsilon_0)=0\).

Conversely, if \(c(\epsilon_0)=0\), there is a section
\[
  w\in H^0(S,j_*L_0)
\]
such that, for every \(S\)-prescheme \(Y\) and every \(g\in G(Y)\),
\[
  z(Y)(g)=g\cdot w_Y-w_Y.
\]
Apply this with \(Y=G\) and \(g=g^\sharp\). It follows that
\(a-w\in X^G(S)\). Thus \(\epsilon_0\) lifts to \(X^G(S)\) if and
only if \(c(\epsilon_0)=0\).

\paragraph{\textbf{Remark 9.5.}}
\SGARefTarget{sga3:xii:9.5}
When \(S\) is affine and \(G\) is affine, flat, and of finite type
over \(S\), this obstruction can be refined to a class
\[
  \widetilde c(\epsilon_0)\in H^1_{G_0}(X_0,L_0),
\]
where the right-hand side is the \(G_0\)-equivariant cohomology
group defined by Wewers (Appendix~C). In this case it is unnecessary
to assume that \(\epsilon_0\) lifts to \(X(S)\).

\setcounter{subsection}{5}
\subsection{Smoothness of fixed points}
\SGARefTarget{sga3:xii:9.6}

\paragraph{\textbf{Theorem 9.7.}}
\SGARefTarget{sga3:xii:9.7}
\emph{Let \(G\) be a flat \(S\)-group scheme of finite type over a
noetherian base \(S\), acting on a smooth \(S\)-scheme \(X\).
Suppose that \(G\) is \(S\)-pure and that, for every geometric point
\(s\) of \(S\) with algebraically closed residue field,}
\[
  H^1\bigl(G_s,(\Omega^1_{X_s})^*\bigr)=0.
\]
\emph{Then the fixed-point functor \(X^G\) is represented by a
closed subscheme of \(X\), smooth over \(S\).}

By 9.2, \(X^G\) is represented by a closed subscheme of \(X\), of
finite presentation over \(S\). We may suppose
\(S=\operatorname{Spec}A\) with \(A\) local. By the smoothness
criterion of SGA~1, III~3.1(iii bis), it is enough to verify formal
smoothness for
\[
  S'=\operatorname{Spec}A'
\]
and its closed subscheme
\[
  S''=\operatorname{Spec}(A'/I'),
\]
where \(A'\) is an artinian local ring and \(I'\) is an ideal of
square zero. By dévissage, we may assume that \(I'\) is annihilated
by the maximal ideal of \(A'\). If \(k\) is the residue field of
\(A'\), then \(I'\) is therefore a \(k\)-vector space.

Let \(\epsilon''\in X^G(S'')\) be a point that we wish to lift to
\(X^G(S')\), and put
\[
\begin{gathered}
  X'=X\times_S S',\qquad G'=G\times_S S',\\
  X''=X\times_S S'',\qquad G''=G\times_S S''.
\end{gathered}
\]
Since \(S'\) is affine and \(X'\) is smooth over \(S'\),
\(\epsilon''\) lifts to \(X(S')\). By \SGARefLink{sga3:xii:9.4}{Lemma~9.4}, the obstruction to
lifting \(\epsilon''\) to \(X^G(S')\) is a class
\[
  c(\epsilon'')\in H^1(G'',L''),
\]
where
\[
  L''=
  \operatorname{Hom}_{\mathcal O_{S''}}
  \bigl((\epsilon'')^*\Omega^1_{X''/S''},I'\bigr).
\]
Because \(L''\) is a \(k\)-vector space, there is a canonical
isomorphism
\[
  H^1(G'',L'')
  \simeq
  H^1(G''\times_{A'/I'}k,L'').
\]
Let \(\bar k\) be an algebraic closure of \(k\). By \SGARefLink{sga3:ix:theorem:3:1}{IX~3.1} in the
affine case, and by \SGARefLink{sga3:xii:9.11}{Lemma~9.11} in general, there is an isomorphism
\[
  H^1(G''\times_{A'/I'}k,L'')\otimes_k\bar k
  \xrightarrow{\ \sim\ }
  H^1(G''\times_{A'/I'}\bar k,L''\otimes_k\bar k).
\]
The representation \(L''\otimes_k\bar k\) is a direct sum of copies
of
\[
  (\Omega^1_{X_{\bar k}})^*.
\]
The group on the right therefore vanishes by hypothesis. It follows
that \(H^1(G'',L'')=0\), so the obstruction vanishes and \(X^G\) is
formally smooth.

\paragraph{\textbf{Corollary 9.8.}}
\SGARefTarget{sga3:xii:9.8}
\emph{Let \(G\) be a flat \(S\)-group scheme of finite presentation
acting on a smooth \(S\)-scheme \(X\) of finite presentation.
Suppose that \(G\) admits a composition series whose factors are of
one of the following types:}
\begin{enumerate}
\item[\emph{(1)}] \emph{an abelian \(S\)-scheme, that is, a smooth
\(S\)-group scheme whose fibers are abelian varieties;}
\item[\emph{(2)}] \emph{an \(S\)-group scheme of multiplicative
type;}
\item[\emph{(3)}] \emph{a finite étale \(S\)-group scheme of degree
invertible on \(S\);}
\item[\emph{(4)}] \emph{a reductive \(S\)-group scheme, when \(S\)
is a \(\mathbb Q\)-scheme.}
\end{enumerate}
\emph{Then the fixed-point functor \(X^G\) is represented by a
closed subscheme of \(X\), of finite presentation and smooth over
\(S\).}

Indeed, suppose that
\[
  1\longrightarrow G'\longrightarrow G\longrightarrow G''
  \longrightarrow 1
\]
is an exact sequence of \(S\)-groups satisfying the hypotheses.
The group \(G''\) acts on \(X^{G'}\), and \(X^G\) is the fixed-point
functor of \(X^{G'}\) under \(G''\). We are therefore reduced to the
case of a group of one of the four types listed above. By passage to
the limit, we may suppose \(S\) noetherian. We verify the
cohomological criterion of \SGARefLink{sga3:xii:9.7}{Theorem~9.7} for a geometric fiber \(G_s\)
and a finite-dimensional linear representation \(V\) of \(G_s\).

For an abelian scheme, every morphism from \(G_s^n\) to \(V\) is
constant; consequently
\[
  H^i(G_s,V)=0\qquad(i>0).
\]
For a group scheme of multiplicative type, \(G_s\) is
diagonalizable, and \SGARefLink{sga3:i:theorem:5-3-3}{I~5.3.3} again gives
\[
  H^i(G_s,V)=0\qquad(i>0).
\]

In the remaining two cases the assertion follows from semisimplicity
of the linear representations of \(G_s\). If \(G_s\) is finite
étale of degree invertible in \(\kappa(s)\), it is a constant finite
group of invertible order, and every linear representation of
\(G_s\) is semisimple. In the reductive case,
\(\kappa(s)\) is algebraically closed of characteristic zero, and
every linear representation of \(G_s\) is semisimple
(Tauvel--Yu, Theorem~27.3.3).

\paragraph{\textbf{Remark 9.9.}}
\SGARefTarget{sga3:xii:9.9}
This applies in particular to an action of \(G/S\) on a smooth
\(S\)-group scheme \(H\). If \(G\) is of multiplicative type and
acts by inner automorphisms on \(H\), with \(H\) smooth and affine
over \(S\), one recovers \SGARefLink{sga3:xi:corollary:5:3}{XI~5.3}: the centralizer of \(G\) in \(H\)
is smooth.

\paragraph{\textbf{Remark 9.10.}}
\SGARefTarget{sga3:xii:9.10}
When \(S=\operatorname{Spec}k\) and \(G\) is a linearly reductive
algebraic group over an algebraically closed field \(k\), the result
is due independently to Fogarty (Theorem~5.4) and Iversen
(Proposition~1.3).

The proof of \SGARefLink{sga3:xii:9.7}{Theorem~9.7} used the following lemma, which is well
known for affine group schemes (\SGARefLink{sga3:ix:theorem:3:1}{IX~3.1}).

\paragraph{\textbf{Lemma 9.11.}}
\SGARefTarget{sga3:xii:9.11}
\emph{Let \(G\) be a group scheme over an affine scheme
\(S=\operatorname{Spec}A\), and let
\(f:S'=\operatorname{Spec}A'\to S\) be flat. Put
\(G'=G\times_S S'\). Then, for every quasi-coherent
\(G\)-\(\mathcal O_S\)-module \(M\), there are canonical
isomorphisms}
\[
  H^i(G,M)\otimes_A A'
  \xrightarrow{\ \sim\ }
  H^i(G',M\otimes_A A')
  \qquad(i\geq 0).
\]

For \(n\geq1\), the degree-\(n\) Hochschild cochains are determined
by their value at
\[
  (\operatorname{id},\ldots,\operatorname{id})
  \in G(G)\times\cdots\times G(G).
\]
Thus there is an isomorphism of \(A\)-modules
\[
  C^n(G,M)
  \simeq
  \Gamma\bigl(G^n,M\otimes_A\mathcal O_{G^n}\bigr).
\]
The group \(H^n(G,M)\) is the \(n\)-th cohomology group of the
complex
\[
  M\longrightarrow
  \Gamma(G,M\otimes_A\mathcal O_G)\longrightarrow
  \Gamma(G^2,M\otimes_A\mathcal O_{G^2})\longrightarrow
  \Gamma(G^3,M\otimes_A\mathcal O_{G^3})\longrightarrow\cdots.
\]
Since \(A'\) is flat over \(A\), EGA IV\(_1\), 1.7.21 gives
\[
\begin{split}
  \Gamma\bigl(G^n,M\otimes_A\mathcal O_{G^n}\bigr)\otimes_A A'
  \xrightarrow{\ \sim\ }{}&
  \Gamma\bigl((G')^n,\\
  &(M\otimes_A A')\otimes_{A'}
  \mathcal O_{(G')^n}\bigr).
\end{split}
\]
Taking cohomology proves the assertion.

\begin{thebibliography}{Ray72}
\SGARefTarget{sga3:xii:bibliography}
\addcontentsline{toc}{section}{Bibliography}

\item[]\emph{Editor's note: the following are additional
references cited in this Exposé.}

\bibitem[Fo73]{sga3-xii-Fo73}
J.\ Fogarty,
\emph{Fixed point schemes},
Amer.\ J.\ Math.\ \textbf{95} (1973), 35--51.

\bibitem[RG71]{sga3-xii-RG71}
M.\ Raynaud and L.\ Gruson,
\emph{Critères de platitude et de projectivité},
Invent.\ Math.\ \textbf{13} (1971), 1--89.

\bibitem[Iv72]{sga3-xii-Iv72}
B.\ Iversen,
\emph{A fixed point formula for action of tori on algebraic
varieties},
Invent.\ Math.\ \textbf{16} (1972), 229--236.

\bibitem[Ray72]{sga3-xii-Ray72}
M.\ Raynaud,
\emph{Flat modules in algebraic geometry},
Compositio Math.\ \textbf{24} (1972), 11--31.

\bibitem[TY06]{sga3-xii-TY06}
P.\ Tauvel and R.\ W.\ T.\ Yu,
\emph{Lie algebras and algebraic groups},
Springer-Verlag, 2006.

\bibitem[We05]{sga3-xii-We05}
S.\ Wewers,
\emph{Formal deformation of curves with group scheme action},
Ann.\ Inst.\ Fourier (Grenoble) \textbf{55} (2005), 1105--1165.

\end{thebibliography}
